% Generated draft: requires mathematical and source review.
\documentclass[11pt]{article}
\usepackage[margin=25mm]{geometry}
\usepackage{fontspec}
\setmainfont{Latin Modern Roman}
\usepackage{amsmath,amssymb,mathtools,unicode-math}
\setmathfont{Latin Modern Math}
\usepackage{enumitem,needspace,xurl,hyperref}
\setlength{\parindent}{0pt}
\setlength{\parskip}{5pt}
\setlength{\emergencystretch}{3em}
\newcommand{\N}{\mathbb{N}}
\newcommand{\Z}{\mathbb{Z}}
\newcommand{\Q}{\mathbb{Q}}
\newcommand{\R}{\mathbb{R}}
\newcommand{\C}{\mathbb{C}}
\newcommand{\F}{\mathbb{F}}
\newcommand{\A}{\mathbb{A}}
\newcommand{\PP}{\mathbb{P}}
\newcommand{\E}{\mathbb{E}}
\newcommand{\eps}{\varepsilon}
\newcommand{\dd}{\,\mathrm{d}}
\begin{document}
\section*{UnsolvedMath ID: 30002697}
\textbf{Source problem code:} catalogue entry
\subsection{Definitions and mathematical statement}
For a real Lie group \(S\) with Lie algebra \(\mathfrak{s}\), define its derived series by
\[
\mathfrak{s}^{(0)}=\mathfrak{s},\qquad
\mathfrak{s}^{(j+1)}=[\mathfrak{s}^{(j)},\mathfrak{s}^{(j)}]\quad(j\geq 0).
\]
The group is called solvable when \(\mathfrak{s}^{(m)}=0\) for some integer \(m\geq 0\). For a Riemannian metric \(g\) on \(S\), its full isometry group is
\[
\operatorname{Isom}(S,g)=\{\phi:S\to S\mid \phi\text{ is a smooth diffeomorphism and }\phi^*g=g\}.
\]
The metric is Einstein if \(\operatorname{Ric}_g=\kappa g\) for some constant \(\kappa\in\R\). The Ricci soliton equation is
\[
\operatorname{Ric}_g+\tfrac12\mathcal{L}_Xg=\lambda g,
\]
where \(X\) is a smooth vector field on \(S\), \(\lambda\in\R\) is constant, and \(\mathcal{L}_X\) denotes Lie differentiation. In this sign convention, expanding solitons have \(\lambda<0\). Completeness requirements for the metric and the soliton vector field are not specified in the supplied record.

\textbf{Source question.} To what extent do maximal-isometry-group results for Einstein metrics extend to Ricci soliton metrics on general solvable Lie groups?

\textbf{Unresolved formulation.} The record does not identify the Einstein theorem or the precise maximality property to be extended. It also does not specify the comparison class of metrics, whether one compares full isometry groups or their identity components, or whether maximality means group inclusion after a change of coordinates, maximal dimension, or another property. Consequently, the supplied question is a research program rather than a uniquely determined quantified conjecture.

No restriction to left-invariant metrics is imposed by the wording supplied. If such a restriction is intended, it means \(L_s^*g=g\) for every \(s\in S\), where \(L_s(t)=st\). Whether this restriction is intended, whether \(S\) must be connected or simply connected, and whether the target concerns all solitons or only expanding solitons require clarification. These hypotheses must be specified before a precise maximal-isometry assertion can be stated.
\subsection{Short English statement}
Determine which maximal-isometry results for Einstein metrics remain valid for Ricci soliton metrics on solvable Lie groups. The supplied record does not identify the comparison theorem, the meaning of maximality, or all required hypotheses, so it does not determine a single precise conjecture.
\subsection{Sources}
\begin{itemize}
\item Local UnsolvedMath registry, permanent ID \(30002697\), source problem code \texttt{catalogue entry}, title \emph{Maximal Isometry Groups of Solvable Ricci Solitons}, in the Oberwolfach Reports Open Problems collection. The record gives proposed year 2014 and names no proposer. Its supplied status is \texttt{partially\_solved}; its \texttt{corrected\_verified} statement label records a supplied semantic, typographic, and LaTeX review, not independent verification here.
\item Supplied report attribution: \emph{Mini-Workshop: Einstein Metrics, Ricci Solitons and Ricci Flow under Symmetry Assumptions} (2015), Oberwolfach Reports, pp.~2529--2568, DOI \texttt{10.4171/owr/2014/45}. \url{https://doi.org/10.4171/owr/2014/45}. The background describes Oberwolfach Reports as documenting workshops and their open-problem discussions. The citation's 2015 date coexists with the record's proposed year 2014 and its designation as Report 45/2014; this discrepancy is unresolved.
\item Supplied literature source: \emph{Every Semialgebraic Ricci Soliton is Algebraic}. \url{https://msp.org/gt/2014/18-4/gt-v18-n4-p12-s.pdf}.
\item Supplied literature source: \emph{Infinitesimal Maximal Symmetry and Ricci Solitons}. \url{https://arxiv.org/abs/2112.09733}.
\item Supplied literature source: \emph{Oberwolfach Report 45/2014}. \url{https://doi.org/10.4171/owr/2014/45}.
\item Literature assessment dated exactly \texttt{2026-08-21}: supplied dated curation, not independently verified. It labels the problem partially solved and reports algebraicity for homogeneous Ricci solitons relative to a suitable transitive isometry group, strong maximal-symmetry and uniqueness results for expanding solvable solitons, and a lack of complete uniformity for the full analogue of the Einstein maximal-isometry theorem. The background explicitly calls this dated literature triage derived from the supplied OWR curation and asks that time-sensitive status and cited claims be verified independently.
\end{itemize}
\end{document}
